
\section{Transgression algebra $B_\Theta$ and secondary anomaly $u=\eta^2$:
the anomaly-completed holographic datum}
\label{app:anomaly-completed-extensions-comparisons}

\providecommand{\ModEnv}{\operatorname{ModEnv}}
\providecommand{\Mbar}{\overline{\mathcal M}}
\providecommand{\Rring}{R^{\bullet}}

\begin{theorem}[Tautological transgression action]
\label{thm:tholog-tautological-transgression-action}
Suppose the modular envelope \(\ModEnv(\cT;B)\) exists and carries
state spaces \(\mathcal{H}_{g,n}^{\cT;B}\), modular operations
\(\mu_{g,n}^{\cT}\), and a universal correlator class
\[
\mathcal{U}_{g,n}^{\cT;B}
\in
H^{\bullet}\!\bigl(\Mbar_{g,n};
\End(\mathcal{H}_{g,n}^{\cT;B})\bigr)
\]
produced by factorization-homology pushforward
(Definition~\ref{def:descendant-taut-action}).  Then the
tautological ring acts by
\[
\Theta_{g,n}^{\cT;B}
\;:\;
\Rring(\Mbar_{g,n})
\;\longrightarrow\;
\End(\mathcal{H}_{g,n}^{\cT;B}),
\qquad
\alpha\;\longmapsto\;
\int\alpha\cup\mathcal{U}_{g,n}^{\cT;B},
\]
and every tautological relation on \(\Mbar_{g,n}\) maps to an operator
identity on \(\mathcal{H}_{g,n}^{\cT;B}\).  The handle part factors
through \(\mathfrak{G}_{g}(B_{\Theta})\) when the handle
representatives are the operators \(\alpha_i,\beta_i\) and the
secondary charge \(u=\eta^2\).
\end{theorem}

\begin{proof}
The displayed assignment is cup product with
\(\mathcal U_{g,n}^{\cT;B}\) followed by integration over
\(\Mbar_{g,n}\).  Multiplicativity is the modular-envelope
composition law.  The handle factorization is the universal property
of Definition~\ref{def:tholog-genus-clifford}.
\end{proof}

\subsection{The anomalous Steinberg}
\label{subsec:anomalous-steinberg}

\providecommand{\Mvac}{\mathcal{M}_{\mathrm{vac}}}
\providecommand{\Lb}{\mathcal{L}_b}
\providecommand{\Sb}{\mathfrak{S}_b}

The Chriss--Ginzburg presentation of the Hecke algebra via the
Steinberg variety $\Sb = \Lb \times_{\Mvac} \Lb$ has a curved
analogue when the vacuum stack carries a pre-symplectic form whose
failure of closedness is the anomaly class.

\begin{construction}[The anomalous Steinberg]
\label{constr:anomalous-steinberg}
\index{Steinberg variety!anomalous}
\index{anomaly completion!Steinberg interpretation}
Let $\cT$ be a 3d HT theory with vacuum moduli stack
$\Mvac(\cT)$ carrying a degree \(-2\) pre-symplectic form
$\omega$ satisfying
\[
d\omega \;=\; H_3,
\]
where $H_3$ is the 't~Hooft anomaly $3$-form.  The Steinberg
object $\Sb = \Lb \times_{\Mvac(\cT)} \Lb$ inherits a
\emph{curved} $(-1)$-shifted structure: the self-intersection
carries a curvature form proportional to $H_3|_{\Lb}$. The
convolution algebra on $\Sb$ becomes curved:
\[
d^2_{\mathrm{conv}} \;=\; [H_3,\;\cdot\;] \;\neq\; 0.
\]
The curved convolution algebra on $\Sb$ is the
anomaly-obstructed algebra of boundary operators.
\end{construction}

\begin{proposition}[Gerbe trivialization criterion]
\label{prop:anomaly-completion-trivialization}
\index{anomaly completion!trivialization on extended stack}
Assume there is an extended stack
\(\widetilde{\Mvac}\to\Mvac\) and a form \(\beta\) on
\(\widetilde{\Mvac}\) whose differential is the pullback of
\(H_3\).  The extended Steinberg
\[
\widetilde{\Sb}
\;=\;
\widetilde{\Lb}
\;\times_{\widetilde{\Mvac}}\;
\widetilde{\Lb}
\]
then carries the uncurved convolution algebra obtained from the
curved convolution algebra on \(\Sb\) by adjoining the trivializing
cochain.  Its modules are boundary operators equipped with a chosen
neutralization of \(H_3\).
\end{proposition}

\begin{proof}
Pulling back the curved differential gives
\(d_{\mathrm{conv}}^2=[d\beta,\cdot]\).  Adding the transgression
cochain with differential \(d\beta\) replaces this curved
differential by an honest differential, exactly as in
Definition~\ref{def:tholog-transgression-algebra}.
\end{proof}

\begin{remark}[Modular MC element on the curved Steinberg]
\label{rem:modular-mc-curved-steinberg}
The modular Maurer--Cartan element of the anomaly-completed datum
lives in the curved convolution algebra
$\mathrm{MC}(\mathfrak{g}^{\mathrm{SC},\mathrm{curved}}_\cT)$
on the anomalous Steinberg $\Sb$. The anomaly completion
lifts this element to the uncurved algebra
$\mathrm{MC}(\mathfrak{g}^{\mathrm{SC}}_{\widetilde\cT})$
on the extended Steinberg $\widetilde{\Sb}$. The passage
from the curved to the uncurved MC equation is precisely the
trivialization of $H_3$: the obstruction $d^2 = [H_3, \cdot]$
is absorbed into the extended stack, and the lifted MC element
satisfies the standard (uncurved) MC equation
$d\alpha + \tfrac{1}{2}[\alpha,\alpha] = 0$.
\end{remark}

\subsection{Abelian Chern--Simons: $B_\Theta = \Bbbk[x,\eta]$ with
$d\eta = kx$ and secondary class $u = \eta^2$}
\label{subsec:abelian-cs-example}

\begin{computation}[Genus-zero holographic datum for abelian CS]
\label{comp:abelian-cs-strict-datum}
\index{Chern--Simons!abelian!genus-zero holographic datum}
For $U(1)$ Chern--Simons at level~$k$ on
$\mathbb R_+\times\mathbb C$:
\begin{enumerate}[label=\textup{(\roman*)}]
\item The boundary algebra is the Heisenberg vertex algebra
 $A_\partial=\cH_k$.
\item Under Remark~\ref{rem:tholog-hmkd-comparison}, the strict
 genus-zero holographic modular Koszul datum is
 \[
 B=\cA^!_{\mathrm{line}},
 \qquad
 A_\partial=\cA=\cH_k,
 \qquad
 \Theta=\frac{k}{z},
 \]
 with the genus-zero shadow connection
 $\nabla^{\mathrm{hol}}=d$.
\item This strict package is \emph{not} yet the anomaly-completed
 datum. The naive odd-generator model
 $B=\Bbbk[\alpha]$ with $|\alpha|=1$ forces
 $\alpha^2=0$, so it cannot encode the genus-one anomaly class.
\item The correct anomaly-completed input is therefore the cochain
 model
 \[
 B=C^\bullet(\mathfrak a)=\Bbbk[x],
 \qquad
 |x|=2,
 \qquad
 \Theta=kx,
 \]
 which is worked out explicitly in
 Example~\ref{ex:abelian-cs-anomaly-completion}.
\end{enumerate}
\end{computation}

\begin{example}[Abelian Chern--Simons: the full anomaly-completion construction]
\label{ex:abelian-cs-anomaly-completion}
\index{Chern--Simons!abelian!anomaly completion}
\index{transgression algebra!abelian Chern--Simons}
\index{genus-Clifford completion!abelian Chern--Simons}
We take the simplest nontrivial input, $U(1)$ Chern--Simons at
level~$k\neq 0$, and trace every abstract construction of this chapter
in fully explicit form. This is the abelian atom of the anomaly-completed
holographic construction.

\medskip
\noindent\textbf{Step 1: the input algebra and anomaly class.}
Let
\[
B \;=\; \Bbbk[\alpha],
\qquad |\alpha|=1,\qquad d=0,
\]
the polynomial algebra on one odd generator (equivalently, the
exterior algebra $\Lambda(\alpha)$, but over a field of
characteristic~$0$ with the Ore/skew convention of
Definition~\ref{def:tholog-transgression-algebra}, we keep
the polynomial notation to match the general formalism).
Because $|\alpha|=1$ is odd,
\[
\alpha^2 \;=\; -\alpha^2
\]
in a supercommutative algebra, so $\alpha^2=0$ and $B\cong\Lambda(\alpha)
= \Bbbk\oplus\Bbbk\alpha$ with $B^0=\Bbbk$, $B^1=\Bbbk\alpha$.

The anomaly class is
\[
\Theta \;:=\; k\,\alpha^2 \;=\; 0 \;\in\; B.
\]
Since $\alpha$ is odd and $B$ is the exterior algebra on one generator,
$\alpha^2=0$ in $B$, so $\Theta=0$.
This is the abelian miracle: the bulk Chern--Simons class
$\frac{k}{2}\int\alpha\,d\alpha$ is quadratic in the gauge field, but
the degree-$2$ element $\alpha^2$ vanishes identically in the exterior
algebra. The anomaly is \emph{zero at the strict polynomial level}.

\medskip
\noindent\textbf{Step 1$'$: the correct input (cochain model).}
The vanishing above shows that the naive exterior-algebra model is too
small. The physically correct bulk algebra uses the \emph{Chevalley--Eilenberg
cochains} of the abelian Lie algebra $\mathfrak{a}=\Bbbk$:
\[
B \;=\; C^\bullet(\mathfrak{a}) \;=\; \Bbbk[x],
\qquad |x|=2,\qquad d=0.
\]
Here $x$ is the degree-$2$ generator corresponding to the curvature
$F=d\alpha$ (the Chern--Weil image of the generator of $\mathfrak{a}^*$).
In this model,
\[
\Theta \;:=\; k\,x \;\in\; B^2 \;=\; \Bbbk\cdot x.
\]
Since $B=\Bbbk[x]$ is a polynomial algebra in an even-degree generator,
$\Theta$ is manifestly closed ($d=0$) and central (every element of
$\Bbbk[x]$ commutes with $x$).

This is the model we use for the remainder of this example.

\medskip
\noindent\textbf{Step 2: the transgression algebra $B_\Theta$ explicitly.}
By Definition~\ref{def:tholog-transgression-algebra},
\[
B_\Theta
\;=\;
\frac{\Bbbk[x]*\Bbbk\langle\eta\rangle}
{\langle\eta\,b-(-1)^{|b|}b\,\eta\;\mid\;b\in B\rangle},
\qquad |\eta|=1,\qquad d\eta=kx.
\]
Since $|x|=2$ (even), the skew-commutation relation gives
$\eta\,x=(-1)^{|x|}x\,\eta = x\,\eta$.
So $\eta$ commutes with $x$, and the transgression algebra is
\[
B_\Theta \;=\; \Bbbk[x,\eta],
\qquad |x|=2,\;\;|\eta|=1,\qquad
\eta x=x\eta,\qquad
d\eta=kx,\;\;dx=0.
\]
This is a polynomial algebra in two generators: one even ($x$, degree $2$)
and one odd ($\eta$, degree $1$), with a nontrivial differential.

By Remark~\ref{rem:tholog-normal-form}, every element has a unique
normal form $\sum_{n\geq 0} p_n(x)\,\eta^n$ with $p_n\in\Bbbk[x]$.
As a graded $\Bbbk[x]$-module, $B_\Theta$ is free on
$\{1,\eta,\eta^2,\eta^3,\dots\}$.

\medskip
\noindent\textbf{Step 3: cohomology of $B_\Theta$.}
The differential acts by $d(\eta)=kx$, $d(x)=0$. On monomials:
\[
d(x^a\eta^b) \;=\; b\,k\,x^{a+1}\eta^{b-1}
\qquad (b\geq 1).
\]
Set $u:=\eta^2\in B_\Theta^2$. By
Proposition~\ref{prop:tholog-secondary-central-class},
$u$ is closed and central. Check directly:
\[
d(u)=d(\eta^2)=(d\eta)\eta-\eta(d\eta)
= kx\eta-\eta\cdot kx=kx\eta-kx\eta=0.
\]
(Here $\eta x=x\eta$ was used.) By
Theorem~\ref{thm:tholog-hypersurface}, $C=B[u]=\Bbbk[x,u]$
and
\[
H^\bullet(B_\Theta) \;\cong\; \Bbbk[x,u]/(kx) \;\cong\; \Bbbk[u],
\]
since multiplication by $[\Theta]=kx$ is injective on $\Bbbk[x]$ (as
$k\neq 0$), so the hypersurface formula applies and the quotient
by the ideal $(kx)$ kills all powers of~$x$.

The cohomology is a polynomial ring in one even variable~$u$ of
degree~$2$:
\[
H^\bullet(B_\Theta) \;\cong\; \Bbbk[u].
\]

\medskip
\noindent\textbf{Step 4: the one-sided resolution.}
Let $M$ be a left dg $B_\Theta$-module. The neutralisation is the
action of $\eta$:
\[
h_M \;=\; \eta_M \;:\; M\to M,\qquad h_M(m):=\eta\cdot m.
\]
This satisfies $d(h_M)=\Theta_M$ (the action of $kx$ on $M$)
because $d\eta=kx$.

By Theorem~\ref{thm:tholog-free-two-term-resolution}, the free
two-term resolution is
\[
0 \;\longrightarrow\;
(B_\Theta\otimes_B M)[-1]
\;\xrightarrow{\;\widetilde D_M\;}
B_\Theta\otimes_B M
\;\xrightarrow{\;\mu_M\;}
M
\;\longrightarrow\; 0,
\]
with $D_M(s\otimes m)=s\eta\otimes m - (-1)^{|s|}s\otimes\eta m$.
Since $B=\Bbbk[x]$ and $B_\Theta=\Bbbk[x,\eta]$, the induced
module $B_\Theta\otimes_B M\cong\Bbbk[\eta]\otimes_\Bbbk M$
as a graded vector space, and $D_M$ acts by
\[
D_M(\eta^n\otimes m) \;=\; \eta^{n+1}\otimes m
\;-\; (-1)^n\,\eta^n\otimes\eta m.
\]

For the trivial module $M=\Bbbk$ (with $x$ and $\eta$ acting by
zero), this becomes $D_\Bbbk(\eta^n\otimes 1)=\eta^{n+1}\otimes 1$,
i.e.\ the resolution is the cone on multiplication by~$\eta$.

\medskip
\noindent\textbf{Step 5: the secondary anomaly $u=\eta^2$ on modules.}
For a $B_\Theta$-module $M$ with neutralisation $h_M=\eta_M$,
the secondary anomaly (Definition~\ref{def:tholog-secondary-anomaly})
is
\[
\sigma(M,h_M) \;=\; [h_M^2] \;=\; [\eta_M^2] \;=\; [u_M]
\;\in\; \operatorname{Ext}^2_B(M,M).
\]
This is the action of the central element $u=\eta^2$ on $M$.

\emph{Example module.}
Take $M=\Bbbk[x]/(x)=\Bbbk$, the residue field, viewed as a
$B_\Theta$-module via $x\mapsto 0$, $\eta\mapsto 0$. Then
$h_M=0$ and $u_M=0$, so $\sigma(M,0)=0$. The secondary anomaly
vanishes on the trivial module.

\emph{Example module with nontrivial secondary anomaly.}
Take $M=\Bbbk[\eta]\cong\Bbbk[t]$ as a $B_\Theta$-module with
$x\mapsto 0$ and $\eta$ acting by multiplication. Then $u_M$
acts as multiplication by $t^2$, giving $\sigma(M)=[t^2]\neq 0$.

\medskip
\noindent\textbf{Step 6: genus-Clifford completion for $g=1$.}
By Definition~\ref{def:tholog-genus-clifford},
\[
\mathfrak{G}_1(B_\Theta)
\;=\;
\frac{B_\Theta\langle\alpha_1,\beta_1\rangle}
{\langle\alpha_1^2,\;\beta_1^2,\;
\alpha_1\beta_1+\beta_1\alpha_1-u\rangle},
\qquad |\alpha_1|=|\beta_1|=1.
\]
In the abelian CS model, $B_\Theta=\Bbbk[x,\eta]$, so
\[
\mathfrak{G}_1(B_\Theta)
\;=\;
\frac{\Bbbk[x,\eta]\langle\alpha_1,\beta_1\rangle}
{\langle\alpha_1^2,\;\beta_1^2,\;
\alpha_1\beta_1+\beta_1\alpha_1-\eta^2\rangle}.
\]
As a left $B_\Theta$-module, $\mathfrak{G}_1(B_\Theta)$ is free of
rank~$4$ with basis $\{1,\alpha_1,\beta_1,\alpha_1\beta_1\}$.
The Clifford relation is
\[
\alpha_1\beta_1+\beta_1\alpha_1 \;=\; \eta^2.
\]

\medskip
\noindent\textbf{Step 7: Morita triviality after inverting $u$.}
Set $R:=B_\Theta[u^{-1}]=\Bbbk[x,\eta,\eta^{-2}]$.
By Theorem~\ref{thm:tholog-one-handle-matrix},
\[
\mathfrak{G}_1(R)
\;\cong\;
\underline{\operatorname{End}}_R(R\oplus R\varepsilon),
\qquad |\varepsilon|=-1,
\]
via the explicit isomorphism
\[
\alpha_1\;\longmapsto\;u\cdot(\varepsilon\wedge -),
\qquad
\beta_1\;\longmapsto\;\iota_\varepsilon.
\]
Check the Clifford relation:
\[
u(\varepsilon\wedge -)\,\iota_\varepsilon
+\iota_\varepsilon\,u(\varepsilon\wedge -)
\;=\;
u\,\bigl((\varepsilon\wedge-)\iota_\varepsilon
+\iota_\varepsilon(\varepsilon\wedge-)\bigr)
\;=\;u\cdot\operatorname{id}
\;=\;\eta^2\cdot\operatorname{id}.
\]
This confirms $\mathfrak{G}_1(R)\cong\operatorname{Mat}_{2\times 2}(R)$,
so
\[
D\bigl(\mathfrak{G}_1(R)\text{-}\mathbf{mod}\bigr)
\;\simeq\;
D\bigl(R\text{-}\mathbf{mod}\bigr).
\]
The genus-$1$ completion is derived Morita-trivial after inverting
$u=\eta^2$, confirming
Theorem~\ref{thm:tholog-morita-triviality} in this case.

For general genus~$g$, by
Theorem~\ref{thm:tholog-morita-triviality},
\[
\mathfrak{G}_g(R)
\;\cong\;
\operatorname{Mat}_{2^g\times 2^g}(R),
\]
so all higher-genus data is Morita-trivial away from $u=0$.

\medskip
\noindent\textbf{Step 8: the strict locus $u=0$.}
Modulo $u=\eta^2$, by Theorem~\ref{thm:tholog-exterior-degeneration},
\[
\mathfrak{G}_g(B_\Theta)/(\eta^2)
\;\cong\;
\bigl(\Bbbk[x,\eta]/(\eta^2)\bigr)
\otimes
\Lambda(\alpha_1,\beta_1,\dots,\alpha_g,\beta_g).
\]
Since $\Bbbk[x,\eta]/(\eta^2)=\Bbbk[x]\otimes\Lambda(\eta)$
is an exterior-algebra extension of $\Bbbk[x]$, the strict-locus
genus completion is
\[
\Bbbk[x]\otimes\Lambda(\eta,\alpha_1,\beta_1,\dots,\alpha_g,\beta_g),
\]
an exterior algebra on $2g+1$ odd generators over $\Bbbk[x]$.
This is the exterior handle algebra for the abelian
theory on the secondary-anomaly cone.

\medskip
\noindent\textbf{Summary.}
For $U(1)$ Chern--Simons at level $k\neq 0$:
\begin{center}
\renewcommand{\arraystretch}{1.2}
\begin{tabular}{ll}
Input algebra & $B=\Bbbk[x]$, $|x|=2$, $d=0$ \\
Anomaly class & $\Theta=kx\in B^2$ (central, closed) \\
Transgression algebra &
 $B_\Theta=\Bbbk[x,\eta]$, $d\eta=kx$ \\
Cohomology &
 $H^\bullet(B_\Theta)\cong\Bbbk[u]$ \\
Secondary anomaly & $u=\eta^2$, closed and central \\
Genus-$g$ Clifford completion &
 $\mathfrak{G}_g(B_\Theta)$: Clifford on $2g$ generators
 with $\{\alpha_i,\beta_j\}=\delta_{ij}\eta^2$ \\
Morita triviality &
 $\mathfrak{G}_g(B_\Theta)[\eta^{-2}]
 \cong\operatorname{Mat}_{2^g}(B_\Theta[\eta^{-2}])$ \\
Strict locus &
 $\mathfrak{G}_g(B_\Theta)/(\eta^2)
 \cong\Bbbk[x]\otimes\Lambda^{2g+1}$
\end{tabular}
\end{center}
Every abstract construction of
Sections~\ref{subsec:tholog-transgression-algebra}--\ref{subsec:tholog-morita-triviality}
is visible and computable by hand.
\end{example}

\subsection{The Virasoro algebra: anomaly completion
at the shadow obstruction tower}
\label{subsec:tholog-virasoro}

The Virasoro algebra is the simplest $W$-algebra whose
anomaly-completed datum is nontrivial. Unlike the abelian
example, the shadow obstruction tower is \emph{infinite} (type~$M$), the
curvature is a genuine function of the central charge, and
the transgression algebra connects two Virasoro algebras at
complementary central charges.

\begin{computation}[Virasoro anomaly-completed datum]
\label{comp:virasoro-anomaly-completion}
\index{Virasoro algebra!anomaly completion|textbf}
\index{anomaly completion!Virasoro}
Let $\mathrm{Vir}_c$ denote the Virasoro vertex algebra at
central charge~$c$.

\emph{Koszul duality.}
The chiral Koszul dual is
(Vol~I, Theorem~A):
\begin{equation}\label{eq:vir-koszul-dual}
\mathrm{Vir}_c^!
\;\simeq\;
\mathrm{Vir}_{26-c}.
\end{equation}
The duality is \emph{same-family}: both sides are Virasoro
algebras, at complementary central charges summing to~$26$.
The self-dual point is $c=13$, where
$\mathrm{Vir}_{13}^!\simeq\mathrm{Vir}_{13}$.

\emph{The curvature as a function of~$c$.}
The modular Koszul curvature $\kappaChHodge(\mathrm{Vir}_c)\in\mathbb Q(c)$
is computed from the cyclic pairing on the bar complex of
$\mathrm{Vir}_c$ (Vol~I, Theorem~D). The scalar $\kappa$-complementarity constant for the Virasoro
family is $\kappaChHodge(\mathrm{Vir}_c) + \kappaChHodge(\mathrm{Vir}_{26-c})
= 13$ (related to the Vol~I Trinity conductor
$K(\mathrm{Vir}) = c + c^{!} = 26$ via
$\kappa + \kappa^{!} = \varrho \cdot K$ with
$\varrho_{\mathrm{Vir}} = 1/2$; see Vol~I
\cref{V1-chap:universal-conductor}):
\begin{equation}\label{eq:vir-complementarity-vol2}
\kappaChHodge(\mathrm{Vir}_c) + \kappaChHodge(\mathrm{Vir}_{26-c}) = 13.
\end{equation}
Key properties:
\begin{enumerate}[label=\textup{(\roman*)}]
\item $\kappaChHodge(\mathrm{Vir}_{13})=13/2$ (self-duality gives
 symmetric anomaly budget, not vanishing curvature);
\item $\kappaChHodge(\mathrm{Vir}_c)\neq \kappaChHodge(\mathrm{Vir}_{26-c})$
 for $c\neq 13$ generically;
\item the \emph{shifted} curvature
 $\tilde\kappaChHodge(\mathrm{Vir}_c)
 := \kappaChHodge(\mathrm{Vir}_c) - 13/2$ is anti-symmetric under the
 duality involution:
 $\tilde\kappaChHodge(\mathrm{Vir}_{26-c})
 =-\tilde\kappaChHodge(\mathrm{Vir}_c)$,
 and $\tilde\kappaChHodge(\mathrm{Vir}_{13})=0$.
\end{enumerate}
The shifted curvature $\tilde\kappa$ measures the
\emph{distance from self-duality}: it is the scalar obstruction
to extending the genus-$0$ Koszul duality to genus~$1$
beyond the symmetric baseline $K/2$.

\emph{The transgression algebra as bridge.}
Set $B=\bar B^{\mathrm{ch}}(\mathrm{Vir}_c)$
(the bar coalgebra as dg~algebra), and choose a central cocycle
representative of the anomaly class
\(\Theta=\kappaChHodge(\mathrm{Vir}_c)\cdot\omega_1\).  The
transgression algebra is
\[
B_\Theta
\;=\;
\frac{B*\Bbbk\langle\eta\rangle}
{\langle\eta b-(-1)^{|b|}b\eta\rangle},
\qquad
d\eta=\kappaChHodge(\mathrm{Vir}_c)\cdot\omega_1,
\]
Through
the Koszul-dual comparison
\(\mathrm{Vir}_c^!\simeq\mathrm{Vir}_{26-c}\), a
\(B_\Theta\)-module \(M\)
carries simultaneously:
\begin{itemize}
\item a $\mathrm{Vir}_c$-module structure (via the
 $B$-action);
\item a chain-level homotopy $h_M=\eta_M$ relating the
 $\mathrm{Vir}_c$-action to the $\mathrm{Vir}_{26-c}$-action
 (via the bar-cohomology/Koszul structure).
\end{itemize}
The neutralisation condition $dh_M=\Theta_M$ is the statement
that the genus-$1$ curvature obstruction for $M$ can be
absorbed into a homotopy connecting the two central charges.

\emph{The secondary anomaly and the shadow obstruction tower.}
The secondary anomaly $u=\eta^2$ is closed and central in
$B_\Theta$. The Virasoro shadow obstruction tower from Vol~I
(type~$M$, infinite depth) produces nonvanishing shadow
projections $\mathrm{Sh}_r(\Theta_{\mathrm{Vir}})$ at every
degree $r\ge 2$: the curvature~$\kappaChHodge(\cA)$ at degree~$2$,
the cubic shadow~$C$ at degree~$3$, the quartic resonance
class $Q^{\mathrm{contact}}_{\mathrm{Vir}}
=10/[c(5c{+}22)]$ at degree~$4$, and the quintic obstruction
(forced nonvanishing, degree~$5$).
The transgression algebra records the first central genus class and
the secondary coordinate \(u\).  A lift of the full MC element to a
\(u\)-adically complete transgression algebra is extra data; under
such a lift, the shadow projections are its \(u\)-adic coefficients.

\emph{The Clifford/exterior dichotomy.}
The genus-$g$ Clifford completion
$\mathfrak G_g(B_\Theta)$ bifurcates according to
whether the secondary class \(u=\eta^2\) acts invertibly or
vanishes on the sector under consideration:
\begin{itemize}
\item \emph{Localized regime} ($u \neq 0$ after localization):
 the Clifford algebra $\mathfrak G_g(B_\Theta)$ is
 nondegenerate, hence Morita trivial:
 $\mathfrak G_g[\eta^{-2}]\cong
 \operatorname{Mat}_{2^g}$. This is a Morita equivalence between
 the localized genus-Clifford completion and the localized
 transgression algebra, not an unqualified equivalence of Virasoro
 module categories.
\item \emph{Exterior degeneration} ($u = 0$ on the sector):
 the Clifford relation $\{\alpha_i,\beta_j\} =
 \delta_{ij}\eta^2$ collapses to anticommutativity.
 The genus-$g$ completion degenerates to the exterior
 algebra
 \((B_\Theta/(\eta^2))\otimes
 \Lambda(\alpha_1,\beta_1,\dots,\alpha_g,\beta_g)\).
 Symmetrically, the dual algebra
 $\mathrm{Vir}_{26-c}$ has $\kappaChHodge(\mathrm{Vir}_{26-c}) = (26{-}c)/2$, so
 its intrinsic transgression has zero scalar curvature at $c = 26$.
 In the gravitational context (matter+ghost system),
 it is the effective curvature
 $\kappa_{\mathrm{eff}} = (c{-}26)/2$ that vanishes at
 $c = 26$; this is scalar anomaly cancellation, not a statement about
 the whole modular-invariance problem.
\end{itemize}
Three points in the $c$-line are distinguished:
\begin{enumerate}[label=\textup{(\alph*)}]
\item $c = 0$: the uncurved quadratic self-dual point.
 Here $\kappaChHodge(\mathrm{Vir}_0) = 0$, the bar complex
 is flat, and the intrinsic transgression algebra
 degenerates to exterior.
\item $c = 13$: the Koszul self-dual point.
 Here $\mathrm{Vir}_{13}^! = \mathrm{Vir}_{13}$,
 $\kappaChHodge(\mathrm{Vir}_{13}) = 13/2 \neq 0$, and $u \neq 0$. The Clifford
 algebra is nondegenerate and the genus tower is
 nontrivial at every genus.
\item $c = 26$: the gravitational degeneration point.
 Here $\kappaChHodge(\mathrm{Vir}_{26}) = 13 \neq 0$
 (the intrinsic curvature does not vanish), but
 $\kappa_{\mathrm{eff}} = \kappaChHodge(\mathrm{Vir}_c)
 + \kappa_{\mathrm{ghost}} = 0$. The matter+ghost scalar anomaly
 cancels; any stronger modular statement requires the full
 matter+ghost modular functor.
\end{enumerate}

\emph{Why the Virasoro is the canonical example.}
The Virasoro is the canonical one-parameter example where:
\begin{enumerate}[label=\textup{(\arabic*)}]
\item the Koszul dual is the \emph{same type} of algebra
 (same-family duality);
\item the curvature is a \emph{single rational function}
 of one parameter ($c$);
\item the shadow obstruction tower is \emph{infinite} (type~$M$),
 showing that the transgression algebra carries genuinely
 infinite-depth data;
\item the self-dual point ($c=13$) is an isolated point in
 moduli, not a special sublocus.
\end{enumerate}
Every structural feature of the anomaly-completed datum (the
transgression bridge, the secondary anomaly, the Morita
triviality, the strict-locus degeneration) is visible and
controlled by the single parameter~$c$.
\end{computation}

\subsection{The $\mathfrak{sl}_3$ DS chain:
anomaly-completed holography from the nilpotent poset}
\label{subsec:tholog-ds-chain}

The Virasoro example is the principal reduction of
$\widehat{\mathfrak{sl}}_2$. The first example where the
DS hierarchy and the anomaly-completed datum interact
nontrivially is the $\mathfrak{sl}_3$ chain:
\[
\widehat{\mathfrak{sl}}_{3,k}
\;\xrightarrow{\;\mathrm{DS}_{\min}\;}
\mathcal B^k
\;\xrightarrow{\;\mathrm{DS}_{\mathrm{sub}\to\mathrm{prin}}\;}
\mathcal W_3.
\]
Each arrow is a quantum Drinfeld--Sokolov reduction, and each
vertex carries its own anomaly-completed datum. A DS arrow transports
that datum precisely when the BRST complex carries the bar comparison,
the curvature class, and the handle operators compatibly.

\begin{construction}[The $\mathfrak{sl}_3$ chain as
anomaly-completed diagram]
\label{constr:sl3-anomaly-chain}
\index{Drinfeld--Sokolov reduction!anomaly-completed functor}
\index{sl3 DS chain@$\mathfrak{sl}_3$ DS chain|textbf}
At each vertex of the nilpotent poset of $\mathfrak{sl}_3$
\textup{(}as computed in Volume~I\textup{)}, the
anomaly-completed datum
$(\mathcal A,\;\mathcal A^!,\;\kappa,\;\Theta,\;B_\Theta)$
is:
\begin{center}
\renewcommand{\arraystretch}{1.3}
\begin{tabular}{lccc}
& \textbf{Level~$0$} & \textbf{Level~$1$} & \textbf{Level~$2$} \\
\hline
Orbit & $\{0\}$ & $(2,1)$ & $(3)$ \\[2pt]
$\mathcal A$ &
 $\widehat{\mathfrak{sl}}_{3,k}$ &
 $\mathcal B^k$ &
 $\mathcal W_3^k$ \\[2pt]
$\mathcal A^!$ &
 $\widehat{\mathfrak{sl}}_{3,-k-6}$ &
 $\mathcal B^{-k-6}$ &
 $\mathcal W_3^{-k-6}$ \\[2pt]
$K^c=c+c'$ & $16$ & $196$ & $100$ \\[2pt]
BRST ghost sector & none & minimal & principal \\[2pt]
central midpoint & affine critical & $98$ & $50$ \\[2pt]
Dual orbit & $\{0\}$ & $(2,1)$ & $(3)$ \\[2pt]
Cross-orbit? & no & no (self-$t$) & no
\end{tabular}
\end{center}

\emph{The vertical structure.}
At each level, the anomaly class $\Theta_i=\kappa_i\cdot\omega_1$
is closed and central in the bar dg~algebra $B_i$, and the
transgression algebra $(B_i)_{\Theta_i}$ bridges the
$W$-algebra $\mathcal A_i$ with its Koszul dual
$\mathcal A_i^!$.

\emph{The horizontal structure: DS as functor.}
When the DS reduction $\mathrm{DS}_f\colon\mathcal A\to\mathcal W$
is compatible with the chiral bar construction, the bar complex
comparison has the form
$B_{\mathcal W}\simeq H^0_{\mathrm{BRST}}(B_{\mathcal A}
\otimes B_{\mathrm{ghost}})$,
and the curvature transforms as
\begin{equation}\label{eq:curvature-ds-transform}
\kappaChHodge(\mathcal W)
\;=\;
\kappaChHodge(\mathcal A)
\;-\;
\kappa_{\mathrm{ghost}}(f,k),
\end{equation}
where $\kappa_{\mathrm{ghost}}(f,k)$ is the ghost system's
contribution to the modular characteristic in the chosen DS
normalization.

\emph{Cross-orbit duality from the chain.}
At level~$1$, the partition $(2,1)$ is self-transpose, so
the Koszul dual has the same modular characteristic as
$\mathcal B^{-k-6}$ (same-orbit in the~$\kappa$ sense). The
anomaly-completed datum at level~$1$ is obtained from the
datum at level~$0$ by applying
$\mathrm{DS}_{\min}$: the curvature shifts by
$\kappa_{\mathrm{ghost}}$, the transgression algebra
acquires the ghost sector, and the genus-Clifford completion
inherits a ghost-extended handle algebra. The
\emph{cross-orbit phenomenon} for other nilpotent types arises
when the DS reduction changes the orbit type of the Koszul
dual: the functoriality of DS on anomaly-completed data
then makes the transgression algebra bridge two
\emph{different} $W$-algebra families.
\end{construction}

\begin{theorem}[DS comparison criterion for anomaly-completed data]
\label{thm:ds-anomaly-functoriality}
\index{Drinfeld--Sokolov reduction!functoriality on anomaly completion}
Let $\mathcal A$ be a vertex algebra with anomaly-completed
datum $(B,\Theta,B_\Theta)$ and let $f\in\mathfrak g$ be a
nilpotent element with DS reduction
$\mathcal W=\mathrm{DS}_f(\mathcal A)$.  Assume:
\begin{enumerate}[label=\textup{(\alph*)}]
\item the BRST complex is exact on the objects under consideration;
\item the chiral bar complex satisfies
\[
B_{\mathcal W}\simeq
H^0_{\mathrm{BRST}}(B_{\mathcal A}\otimes\mathcal F_{\mathrm{gh}});
\]
\item the curvature class satisfies
\[
\kappaChHodge(\mathcal W)=\kappaChHodge(\mathcal A)-\kappa_{\mathrm{ghost}}(f,k);
\]
\item the BRST differential commutes with the transgression generator
\(\eta\) and the handle generators \(\alpha_i,\beta_i\).
\end{enumerate}
Then:
\begin{enumerate}[label=\textup{(\roman*)}]
\item The transgression algebra transforms by
 $(B_{\mathcal W})_{\Theta_{\mathcal W}}
 \simeq
 H^0_{\mathrm{BRST}}\bigl(
 (B_{\mathcal A})_{\Theta_{\mathcal A}}
 \otimes\mathcal F_{\mathrm{gh}}
 \bigr)$.
\item The genus-$g$ Clifford completion transforms:
 $\mathfrak G_g(B_{\mathcal W})
 \simeq
 H^0_{\mathrm{BRST}}\bigl(
 \mathfrak G_g(B_{\mathcal A})
 \otimes\mathcal F_{\mathrm{gh}}
 \bigr)$.
\end{enumerate}
In particular, if the anomaly-completed datum at
$\mathcal A$ is Morita-trivial at genus~$g$ (away from the
$u$-locus), then so is the datum at $\mathcal W$.
\end{theorem}

\begin{proof}
By (b) and (c), \(d\eta=\Theta_{\mathcal A}\) descends under BRST
cohomology to \(d\eta=\Theta_{\mathcal W}\).  Since \(Q_{\mathrm{DS}}\)
commutes with \(\eta\) by (d), adjoining \(\eta\) before BRST
cohomology and adjoining it after BRST cohomology give isomorphic
transgression algebras.  The same argument applies to the handle
generators \(\alpha_i,\beta_i\), and their Clifford relations are
preserved because \(u=\eta^2\) is preserved.  Morita triviality after
inverting \(u\) is invariant under the displayed isomorphism.
\end{proof}

\begin{remark}[The $\mathfrak{sl}_3$ chain as the
platonic model]
\label{rem:sl3-platonic-model}
The $\mathfrak{sl}_3$ DS chain is the simplest instance of
the general principle: the nilpotent poset of~$\mathfrak g$
indexes a \emph{diagram} of anomaly-completed data, with the
DS functors as the morphisms. The anomaly class $\Theta_i$ at
each vertex is \emph{not} independent: it is determined by the
curvature transform~\eqref{eq:curvature-ds-transform} from
the affine-level datum. The self-dual locus
varies across the diagram: at level~$0$ it is $k=-3$
(critical; $\kappaChHodge(\widehat{\mathfrak{sl}}_{3,k})=0$), at level~$1$ it is $c=38$,
and at level~$2$ it is $c=50$.
For KM the self-dual and $\kappaChHodge(\cA)=0$ loci coincide;
for $\mathcal{W}$-algebras self-duality means
$\kappaChHodge(\cA) = \kappaChHodge(\cA^!) \ne 0$ (curved self-duality).

The anomaly-completed datum at each vertex is
the \emph{fibre} of this diagram over the specific nilpotent
orbit; the DS functors are the \emph{transition maps} between
fibres.

For $\mathfrak{sl}_N$ with hook partitions, the same
structure governs the entire reduction network $\Gamma_N$
\textup{(}by the transport-propagation proposition of Volume~I\textup{)}:
the transport-propagation lemma says that the
anomaly-completed diagram is \emph{connected}, and the
transport-to-transpose comparison criterion
\textup{(}from Volume~I\textup{)}
is the condition for completeness.
\end{remark}

\begin{remark}[Geometric localization reading of the DS chain]
\label{rem:ds-chain-geometric-reading}
\index{geometric localization!anomaly chain}
\index{Slodowy slice!anomaly chain}
The three vertices of the $\mathfrak{sl}_3$ DS chain admit a
uniform geometric reading through Slodowy slices. At each level
$i\in\{0,1,2\}$, the corresponding nilpotent orbit $\mathcal O_i$
determines a transversal Slodowy slice
$S_{f_i}\subset\mathfrak{g}^*$:
\begin{center}
\renewcommand{\arraystretch}{1.3}
\begin{tabular}{lccc}
& \textbf{Level~$0$} & \textbf{Level~$1$} & \textbf{Level~$2$} \\
\hline
Orbit & $\{0\}$ & $(2,1)$ & $(3)$ \\[2pt]
Slice & $\mathfrak{g}^*$ & $S_{f_{\min}}$ & $S_{f_{\mathrm{prin}}}$
 (Kostant section)
\end{tabular}
\end{center}
The horizontal DS arrows are Hamiltonian reductions that
progressively cut the coadjoint representation from the full
dual~$\mathfrak{g}^*$ to the minimal Slodowy
slice~$S_{f_{\min}}$ to the Kostant
section~$S_{f_{\mathrm{prin}}}$. At each stage the reduction
quotients out by the centraliser of the nilpotent element,
shrinking the transversal slice by exactly the codimension of the
orbit.

The anomaly class at each vertex is the curvature of the arc
space of the corresponding Slodowy slice: the arc space
$J_\infty S_{f_i}$ carries a canonical symplectic form (from
the Kirillov--Kostant--Souriau structure on~$\mathfrak{g}^*$
restricted to~$S_{f_i}$), and the modular
characteristic~$\kappaChHodge(\mathcal W^k(\mathfrak{g},f_i))$ is
the chiral quantisation of its curvature class. The
transformation law~\eqref{eq:curvature-ds-transform},
\[
\kappaChHodge(\mathcal W)
\;=\;
\kappaChHodge(\mathcal A)
\;-\;
\kappa_{\mathrm{ghost}}(f,k),
\]
is the chiral incarnation of the restriction of the
Killing-form curvature to successively smaller slices: the ghost
contribution~$\kappa_{\mathrm{ghost}}(f,k)$ accounts for the
curvature lost in passing from $J_\infty S_{f_i}$ to
$J_\infty S_{f_{i+1}}$. In particular, the level-dependence of
$\kappa_{\mathrm{ghost}}$ at the non-principal step (level~$1$)
reflects the fact that the minimal Slodowy slice is \emph{not} a
linear subspace of~$\mathfrak{g}^*$: the restriction of the
Killing form to~$S_{f_{\min}}$ acquires a level-dependent
correction.

This reading is an instance of the geometric localization
principle developed in Volume~I
: the
shadow obstruction tower at each vertex is the tower of curvature invariants
of the corresponding arc-space geometry, and the DS chain is the
\emph{restriction diagram} of these invariants along the
Slodowy-slice inclusions
$S_{f_{\mathrm{prin}}}\hookrightarrow
S_{f_{\min}}\hookrightarrow\mathfrak{g}^*$.
\end{remark}

\subsection{Hook-type W-algebras: cross-orbit anomaly completion}
\label{subsec:tholog-hook-type}

The abelian Chern--Simons example above is self-dual: the boundary
algebra and its Koszul dual are of the same type. For non-principal
$W$-algebras, the Koszul-dual comparison can cross nilpotent orbits:
the dual algebra is attached to the transpose partition. The
transgression algebra $B_\Theta$ bridges two distinct $W$-algebra
families exactly when the cross-orbit Koszul map identifies their
central genus cocycles.

\begin{construction}[Hook-type anomaly-completed datum]
\label{constr:hook-type-anomaly-datum}
\index{W-algebra@$\mathcal W$-algebra!hook-type!anomaly completion|textbf}
\index{anomaly completion!non-principal W-algebra}
Let $\lambda=(N{-}r,1^r)\vdash N$ be a hook partition of~$N$
with $1\le r\le N{-}2$, and let $\mathcal W_\lambda^k
:= \mathcal W^k(\mathfrak{sl}_N,f_\lambda)$ be the corresponding
$W$-algebra at level~$k\neq -N$. The anomaly-completed datum is the
following sequence of BRST, Koszul-dual, and transgression data.

\emph{Step~1: The BRST input.}
The quantum Drinfeld--Sokolov reduction produces
$\mathcal W_\lambda^k$ as the BRST cohomology
\[
\mathcal W_\lambda^k
\;=\;
H^0\!\bigl(
\widehat{\mathfrak{sl}}_{N,k}
\otimes
\mathcal F_{\mathrm{gh}}(\mathfrak m^+_\lambda),\;
Q_{\mathrm{DS}}^{(\chi_\lambda)}
\bigr),
\]
where $\mathfrak m^+_\lambda$ is the nilradical of the
Jacobson--Morozov parabolic for $f_\lambda$,
$\dim(\mathfrak m^+_\lambda)=r(N{-}r)$, and $\chi_\lambda$
is the Slodowy character. The $W$-algebra has strong generators
of conformal weights $1,2,\dots,r{+}1$ (from the
$\mathfrak{sl}_2$-grading of
$\mathfrak{sl}_N^{f_\lambda}$).

\emph{Step~2: Cross-orbit Koszul duality.}
By the hook-type transport theorem
of Volume~I,
at generic level:
\begin{equation}\label{eq:hook-koszul-duality}
(\mathcal W_\lambda^k)^!
\;\simeq\;
\mathcal W_{\lambda^t}^{k^\vee},
\qquad
k^\vee=-k-2N,
\qquad
\lambda^t=(r{+}1,1^{N{-}r{-}1}).
\end{equation}
The partition transpose $\lambda\mapsto\lambda^t$ is the type-A
specialisation of Barbasch--Vogan duality. For $r\neq(N{-}1)/2$,
the partitions $\lambda\neq\lambda^t$ and the Koszul dual is a
\emph{genuinely different} $W$-algebra.

The central-charge duality identity:
\begin{equation}\label{eq:hook-cc-sum}
c(\mathcal W_\lambda^k)
+c(\mathcal W_{\lambda^t}^{k^\vee})
\;=\;
c_{\mathrm{sum}}(N,r),
\end{equation}
where $c_{\mathrm{sum}}(N,r)$ is the Freudenthal--de~Vries
constant for the pair $(\lambda,\lambda^t)$. For the
Bershadsky--Polyakov algebra ($N=3$, $r=1$):
Volume~I Proposition~\ref*{V1-prop:partition-dependent-complementarity}
proves $c_{\mathrm{sum}}(3,1) = 196$, using the central-charge
formula
\[
c(\mathcal{B}^k)=2-\frac{24(k+1)^2}{k+3}.
\]

\emph{Step~3: The curvature and anomaly class.}
The chiral Koszul curvature $\kappaChHodge(\mathcal W_\lambda^k)$ is
the modular characteristic of $\mathcal W_\lambda^k$,
computed from the bar complex via the cyclic pairing
(Vol~I, Theorem~D). It is a rational function of the
level~$k$.  The shifted curvature
\[
\widetilde\kappaChHodge(\mathcal W_\lambda^k)
:=\kappaChHodge(\mathcal W_\lambda^k)-\frac{1}{2}K^\kappa_\lambda,
\qquad
K^\kappa_\lambda=
\kappaChHodge(\mathcal W_\lambda^k)+
\kappaChHodge(\mathcal W_{\lambda^t}^{k^\vee}),
\]
vanishes at a self-dual point.  The intrinsic curvature
\(\kappaChHodge(\mathcal W_\lambda^k)\) need not vanish there.  For generic
$k$,
$\kappaChHodge(\mathcal W_\lambda^k)\neq 0$ and the genus-$1$ curvature is
$d^2=\kappaChHodge(\mathcal W_\lambda^k)\cdot\omega_1$. The anomaly class is
$\Theta=\kappaChHodge(\mathcal W_\lambda^k)\cdot\omega_1\in Z^2(B)\cap Z(B)$, where $B$
is the bar dg~algebra of the boundary system.

\emph{Step~4: The transgression algebra.}
Set $B=\bar B^{\mathrm{ch}}(\mathcal W_\lambda^k)$
(the bar coalgebra viewed as dg~algebra via the convolution
product) and assume the genus class has been represented by a closed
central cocycle \(\Theta\). The transgression algebra is
\[
B_\Theta
\;=\;
\frac{B*\Bbbk\langle\eta\rangle}
{\langle\eta b-(-1)^{|b|}b\eta\rangle},
\qquad
d\eta=\Theta=\kappa\cdot\omega_1,
\quad|\eta|=1.
\]
The secondary anomaly is
$u=\eta^2\in Z^2(B_\Theta)\cap Z(B_\Theta)$.

The bridge to \(\mathcal W_{\lambda^t}^{k^\vee}\) is the Koszul-dual
comparison map~\eqref{eq:hook-koszul-duality}. A
\(\Theta\)-neutralisation \(h_M\) with \(dh_M=\Theta_M\) on a
\(B\)-module \(M\) is an intertwining datum for that comparison
exactly when the bar cohomology action transported
through~\eqref{eq:hook-koszul-duality} agrees with the \(B\)-action
up to the homotopy \(h_M\).

\emph{Step~5: Genus-Clifford completion.}
The genus-$g$ Clifford completion
\(\mathfrak G_g(B_\Theta)\) has \(2g\) generators
$\alpha_1,\beta_1,\dots,\alpha_g,\beta_g$ with
$\alpha_i\beta_j+\beta_j\alpha_i=\delta_{ij}\,\eta^2$ and
$\alpha_i^2=\beta_j^2=0$.

\begin{enumerate}[label=\textup{(\roman*)}]
\item \emph{Away from the $u$-locus} ($\eta^2$ invertible):
 $\mathfrak G_g(B_\Theta)[\eta^{-2}]
 \cong\operatorname{Mat}_{2^g}(B_\Theta[\eta^{-2}])$
 (Morita triviality).
 Higher genus is Morita-invisible for modules over the localized
 completion. Cross-orbit duality on this locus is controlled by
 whether the Koszul-dual comparison map identifies the localized
 cocycles on the two bar models.

\item \emph{On the strict locus} ($u=\eta^2=0$):
 $\mathfrak G_g(B_\Theta)/(\eta^2)
 \cong (B_\Theta/(\eta^2))\otimes
 \Lambda(\alpha_1,\beta_1,\dots,\alpha_g,\beta_g)$.
 Thus the handle algebra is exterior.  If multiplication by
 \([\Theta]\) is injective in the hypersurface sequence, the cohomology
 of the transgression factor is the corresponding quotient by
 \([\Theta]\); no chain-level quotient \(B/(\Theta)\) is asserted
 without that hypothesis.
\end{enumerate}
\end{construction}

\begin{computation}[Bershadsky--Polyakov algebra:
$\mathfrak{sl}_3$, hook $(2,1)$]
\label{comp:bp-anomaly-completion}
\index{Bershadsky--Polyakov algebra!anomaly completion|textbf}
\index{W-algebra@$\mathcal W$-algebra!subregular!anomaly completion}
The Bershadsky--Polyakov algebra
$\mathcal B^k=\mathcal W^k(\mathfrak{sl}_3,f_{(2,1)})$
is the simplest non-principal case, with $N=3$, $r=1$.

\emph{Generators.}
Four strong generators: a $\widehat{\mathfrak u}_1$ current \(J\)
of conformal weight~\(1\), two fields \(G^\pm\) of conformal
weight~\(3/2\), and the stress tensor \(T\) of conformal weight~\(2\).
This is the standard Bershadsky--Polyakov
\(\mathcal W^{(2)}_3\) generating pattern, not an identification
with the \(\mathcal N=2\) superconformal algebra.

\emph{Koszul dual.}
The partition $(2,1)$ is self-transpose:
$(2,1)^t=(2,1)$. Therefore:
\[
(\mathcal B^k)^!
\;\simeq\;
\mathcal B^{-k-6},
\]
and the Koszul dual is the \emph{same type} of $W$-algebra
at the Feigin--Frenkel dual level $k^\vee=-k-6$.
By Volume~I
Proposition~\ref*{V1-prop:partition-dependent-complementarity},
with \(c(\mathcal{B}^k)=2-24(k+1)^2/(k+3)\):
\[
c(\mathcal B^k)+c(\mathcal B^{-k-6})=196.
\]

\emph{Curvature.}
\[
\kappaChHodge(\mathcal B^k)
\;\in\;\mathbb Q(k),
\]
a rational function of~$k$ computed from the cyclic pairing
on the bar complex of $\mathcal B^k$ (Vol~I, Theorem~D).
Since $(\mathcal B^k)^!=\mathcal B^{-k-6}$ and
$c+c'=196$, the shifted central charge \(c-98\) is odd under
\(k\mapsto -k-6\).  The fixed level is \(k=-3\), the critical
level, so no regular level realizes \(c=98\) in this convention.
The scalar \(\kappa+\kappa^!=98/3\) is the Vol~I
Bershadsky--Polyakov complementarity value.
The explicit formula for $\kappaChHodge(\mathcal B^k)$ is determined
by the bar complex computation
\textup{(}computed in Volume~I\textup{)}:
it receives contributions from both the
$\widehat{\mathfrak{sl}}_2$ sector (double poles) and the
fermionic $G^\pm$ sector (triple pole in the $G^+G^-$ OPE).

\emph{Anomaly class and transgression.}
For a chosen central cocycle representative of the genus class, the
anomaly class is
\(\Theta=\kappa\cdot\omega_1\in Z^2(B)\cap Z(B)\),
where $B=\bar B^{\mathrm{ch}}(\mathcal B^k)$. The
transgression algebra $B_\Theta$ carries the combined
structure of $\mathcal B^k$ and $\mathcal B^{-k-6}$ only after the
same-orbit Koszul comparison has transported the bar action to the
dual level; the neutralisation datum \(h_M\) records the homotopy
of the transported anomaly action.

\emph{Ghost system.}
The DS reduction from $\widehat{\mathfrak{sl}}_3$ to
$\mathcal B^k$ removes $r(N{-}r)=1\cdot 2=2$ positive roots,
producing a ghost contribution
\(\kappa_{\mathrm{ghost}}^{(2,1)}(k)\) in the BRST curvature
formula.  In the Bershadsky--Polyakov normalization used by Vol~I,
this contribution is fixed by the conductor identity
\[
c(\mathcal B^k)+c(\mathcal B^{-k-6})=196,
\qquad
\kappaChHodge(\mathcal B^k)+\kappaChHodge(\mathcal B^{-k-6})=98/3.
\]

\emph{Secondary anomaly and genus structure.}
$u=\eta^2$ is closed and central in $B_\Theta$. The
genus-$1$ Clifford completion
$\mathfrak G_1(B_\Theta)=B_\Theta\langle\alpha_1,\beta_1
\mid\alpha_1\beta_1+\beta_1\alpha_1=\eta^2\rangle$
is the algebraic handle extension at genus \(1\). After inverting
\(\eta^2\), the genus-Clifford completion is Morita equivalent to
the localized transgression algebra. On the strict locus
\(\eta^2=0\), the handle algebra is exterior; comparison with the
dual level still depends on the same-orbit Koszul map.
\end{computation}

\begin{construction}[$\mathfrak{sl}_4$: three-orbit anomaly completion]
\label{constr:sl4-hook-anomaly-completion}
\index{W-algebra@$\mathcal W$-algebra!hook-type!sl4@$\mathfrak{sl}_4$}
\index{anomaly completion!sl4@$\mathfrak{sl}_4$|textbf}
The nilpotent poset of $\mathfrak{sl}_4$ contains five orbits
indexed by partitions of~$4$:
\[
(1^4)\;\subset\;(2,1^2)\;\subset\;(2,2)\;\subset\;(3,1)\;\subset\;(4).
\]
The extreme orbits (zero, principal) are treated by the affine
and principal-$\mathcal W_4$ theory.
The three intermediate orbits ($(3,1)$, $(2,2)$, and
$(2,1,1)$) produce three qualitatively distinct anomaly-completion
patterns, the first of which involves a genuinely cross-orbit
Koszul duality not present for $\mathfrak{sl}_3$.

\emph{Global data.}
$h^\vee=4$, $\dim(\mathfrak{sl}_4)=15$,
$\kappaChHodge(V_k(\mathfrak{sl}_4))=\frac{15(k+4)}{8}$,
$k^\vee=-k-8$.

\smallskip
\emph{Pattern I: the subregular $(3,1)$.}
The subregular $\mathcal W$-algebra
$\mathcal W^k(\mathfrak{sl}_4,f_{(3,1)})$ has \textbf{five strong
generators} of conformal weights $1,2,2,2,3$ (from the
Jacobson--Morozov grading with $\dim=1{+}4{+}5{+}4{+}1$).
The centraliser $\mathfrak{g}^{f_{(3,1)}}$ has dimension~$5$
in JM grades $0,-2,-4$.
Ghost dimension: $r(N{-}r)=1\cdot 3=3$ pairs.
The cross-orbit Koszul duality is
\[
\bigl(\mathcal W^k(\mathfrak{sl}_4,f_{(3,1)})\bigr)^!
\;\simeq\;
\mathcal W^{-k-8}(\mathfrak{sl}_4,f_{(2,1,1)}).
\]
The subregular (5~generators) and minimal (9~generators) $\mathcal W$-algebras
are genuinely distinct: neither the generator count
nor the conformal-weight spectrum is preserved across the duality.

\smallskip
\emph{Pattern II: the rectangular $(2,2)$.}
The partition $(2,2)$ is \emph{not a hook}: it is the first
non-hook, non-principal orbit in this anomaly-completed comparison.
The $\mathcal W$-algebra
$\mathcal W^k(\mathfrak{sl}_4,f_{(2,2)})$ has \textbf{seven strong
generators}: three at weight~$1$ (an $\widehat{\mathfrak{sl}}_2$
subalgebra) and four at weight~$2$.
The JM grading is concentrated in even degrees ($\dim=4{+}7{+}4$).
Since $(2,2)^t=(2,2)$, the Koszul duality is \emph{same-orbit}:
\[
\bigl(\mathcal W^k(\mathfrak{sl}_4,f_{(2,2)})\bigr)^!
\;\simeq\;
\mathcal W^{-k-8}(\mathfrak{sl}_4,f_{(2,2)}).
\]
The self-dual point is $k=-4$ (the critical level).
The complementarity sum $\kappa+\kappa^!$ is a constant
\[
\kappa_0^{(2,2)}
:=
\kappaChHodge(\mathcal W^k(\mathfrak{sl}_4,f_{(2,2)}))
+\kappaChHodge(\mathcal W^{-k-8}(\mathfrak{sl}_4,f_{(2,2)})).
\]
The obstruction to same-orbit anomaly completion is the shifted scalar
\(\kappa-\kappa_0^{(2,2)}/2\); no numerical value for
\(\kappa_0^{(2,2)}\) is used in this comparison.

\smallskip
\emph{Pattern III: the minimal $(2,1,1)$.}
The minimal $\mathcal W$-algebra
$\mathcal W^k(\mathfrak{sl}_4,f_{(2,1,1)})$ has \textbf{nine strong
generators}: $4\times\Delta{=}1$, $4\times\Delta{=}3/2$
(fermionic), $1\times\Delta{=}2$.
Ghost dimension: $r(N{-}r)=2\cdot 2=4$ pairs.
The cross-orbit Koszul duality is
\[
\bigl(\mathcal W^k(\mathfrak{sl}_4,f_{(2,1,1)})\bigr)^!
\;\simeq\;
\mathcal W^{-k-8}(\mathfrak{sl}_4,f_{(3,1)}).
\]

\emph{The $(k+4)$ law.}
At the affine level, $\kappa\propto(k+h^\vee)=(k+4)$ for all
$\mathfrak{sl}_4$ algebras. The orbit-by-orbit curvature is
$\kappaChHodge(\mathcal W_\lambda^k)
=\frac{15(k+4)}{8}-\kappa_{\mathrm{ghost}}^{(\lambda)}(k)$,
where the ghost correction $\kappa_{\mathrm{ghost}}^{(\lambda)}$
is either a constant or a rational function of~$k$ depending on
the parity of the JM grading (see the JM parity criterion below).
\begin{center}
\providecommand{\Ymark}{\text{yes}}
\providecommand{\Nmark}{\text{no}}
\renewcommand{\arraystretch}{1.2}
\begin{tabular}{lccc}
& $(3,1)$ & $(2,2)$ & $(2,1,1)$ \\
\hline
$\dim(\mathfrak{g}^{f_\lambda})$ & $5$ & $7$ & $9$ \\
Generator spectrum &
 $1,2^3,3$ &
 $1^3,2^4$ &
 $1^4,\frac{3}{2}^4,2$ \\
Ghost pairs & $3$ & $4$ & $4$ \\
Dual orbit & $(2,1,1)$ & $(2,2)$ & $(3,1)$ \\
Cross-orbit? & \Ymark & \Nmark\ (self-$t$) & \Ymark \\
$\kappa_{\mathrm{ghost}}$ level-dep? & \Nmark & \Nmark & \Ymark \\
$\kappa\propto(k+4)$? & exact & exact & rational correction
\end{tabular}
\end{center}

\emph{Cross-orbit Koszul duality.}
The pair $(3,1)\leftrightarrow(2,1,1)$ is the $\mathfrak{sl}_4$
incarnation of Barbasch--Vogan duality:
$5$~generators at weights $1,2,2,2,3$ are Koszul dual to
$9$~generators at weights $1^4,(3/2)^4,2$.
The ghost system changes across the duality ($3$ vs $4$ pairs).
The transgression algebra $B_\Theta$ carries
$d\eta=\Theta=\kappa\cdot\omega_1$.  It encodes the genus-$1$
obstruction for both $\mathcal W$-algebras only after the
cross-orbit Koszul comparison identifies their central genus classes.
The neutralisation \(h_M\) is the named homotopy for that comparison.

\emph{The JM parity criterion for level-dependence.}
Orbits with \emph{even-integer} JM grading ($(3,1)$, $(2,2)$)
have level-independent $\kappa_{\mathrm{ghost}}$, because
the ghost conformal weights are fixed integers.
Orbits with \emph{odd-integer} JM grading ($(2,1,1)$, where
the Dynkin good grading introduces half-integer grades)
have $\kappa_{\mathrm{ghost}}(k)$ a nonconstant rational function
of~$k$. This is the \textbf{JM parity criterion}:
\emph{level-dependence of the ghost modular characteristic is
controlled by the parity of the Jacobson--Morozov grading.}

\emph{Genus-Clifford completion.}
For each orbit $\lambda$, the transgression algebra
$B_\Theta^{(\lambda)}$ carries the secondary anomaly
$u=\eta^2\in Z^2(B_\Theta^{(\lambda)})\cap Z(B_\Theta^{(\lambda)})$.
The genus-$g$ Clifford completion
$\mathfrak G_g(B_\Theta^{(\lambda)})$ carries $2g$ handle generators
with $\{\alpha_i,\beta_j\}=\delta_{ij}\eta^2$. Away from $u=0$:
Morita triviality. On the strict locus $u=0$: exterior-algebra
degeneration. For the cross-orbit pair, the strict locus is the
place where the comparison can fail by loss of the transported
central cocycle.
\end{construction}

\begin{remark}[Three duality patterns in $\mathfrak{sl}_4$]
\label{rem:sl4-three-duality-patterns}
\index{anomaly completion!sl4@$\mathfrak{sl}_4$!duality patterns}
The five orbits of $\mathfrak{sl}_4$ organise into three duality patterns:
\begin{enumerate}[label=\textup{(\roman*)}]
\item \emph{Same-family} (affine, principal):
 the affine row has $\kappa+\kappa^!=0$; the principal row has the
 principal \(\mathcal W_4\) complementarity constant determined by
 the Vol~I conductor normalization.
\item \emph{Same-orbit, non-principal} ($(2,2)$):
 the transgression algebra bridges two copies of the same
 $\mathcal W$-algebra family at dual levels, with
 $\kappa+\kappa^!=\kappa_0^{(2,2)}$.
\item \emph{Cross-orbit} ($(3,1)\leftrightarrow(2,1,1)$):
 the transgression algebra bridges two structurally inequivalent
 $\mathcal W$-algebras. This pattern does not exist for $\mathfrak{sl}_3$.
\end{enumerate}
Comparison with $\mathfrak{sl}_3$: for $\mathfrak{sl}_3$, all
non-principal orbits are self-transpose
(the Bershadsky--Polyakov $(2,1)^t=(2,1)$), so only same-orbit
duality appears. The $\mathfrak{sl}_4$ computation is the first
instance of cross-orbit Koszul duality in the anomaly-completed
comparison.
\end{remark}

\begin{remark}[The general hook-type pattern]
\label{rem:hook-anomaly-general}
\index{anomaly completion!hook-type pattern}
For the full hook family
$\mathcal W^k(\mathfrak{sl}_N,f_{(N-r,1^r)})$:
\begin{center}
\renewcommand{\arraystretch}{1.2}
\begin{tabular}{lll}
\textbf{Datum} & \textbf{Formula} & \textbf{Type} \\
\hline
Partition & $\lambda=(N{-}r,1^r)$ & hook \\
Transpose & $\lambda^t=(r{+}1,1^{N{-}r{-}1})$ & hook \\
Self-dual? & iff $N{-}r=r{+}1$, i.e.\ $N=2r{+}1$ (odd) & \\
Dual level & $k^\vee=-k-2N$ & Feigin--Frenkel \\
Ghost dim & $r(N{-}r)$ & nilradical \\
Ghost dim$'$ & $(N{-}r{-}1)(r{+}1)$ & dual nilradical \\
Cross-orbit? & iff $r\neq(N{-}1)/2$ & genuine \\
\end{tabular}
\end{center}
The cross-orbit anomaly completion is the geometrical expression of
the fact that Koszul duality for non-principal
$W$-algebras does not preserve the nilpotent type: the
anomaly $\Theta$ is the obstruction to extending the
genus-$0$ cross-orbit duality to higher genus, and the
transgression algebra $B_\Theta$ is the universal resolution
of this obstruction.
The $\mathfrak{sl}_4$ computation above
(Construction~\ref{constr:sl4-hook-anomaly-completion}) records the
hook pattern and two additional phenomena
not present for hooks alone: the rectangular self-dual orbit $(2,2)$
and the JM parity criterion for level-dependence.
\end{remark}

\subsection{Non-abelian Chern--Simons: the \texorpdfstring{$SU(2)$}{SU(2)}
anomalous Steinberg}
\label{subsec:su2-anomalous-steinberg}

The abelian example of \S\ref{subsec:abelian-cs-example} is the
rank-one commutative model. The first non-abelian case is
\(SU(2)\) Chern--Simons at level~\(k\) on \(\mathbb C\times\mathbb R\).
Every structural feature of the anomaly-completed framework (transgression
algebra, secondary anomaly, anomalous Steinberg, genus-Clifford
completion) acquires genuinely new content from the non-commutativity of
$\mathfrak{sl}_2$.

\begin{construction}[$SU(2)$ Chern--Simons: the anomaly-completed comparison]
\label{constr:su2-cs-anomaly-completion}
\index{Chern--Simons!SU(2)@$SU(2)$!anomaly completion|textbf}
\index{anomaly completion!non-abelian Chern--Simons}
\index{transgression algebra!non-abelian Chern--Simons}
\index{Steinberg variety!anomalous!non-abelian}

\medskip
\noindent\textbf{Step 1: the boundary algebra and its Koszul dual.}
The boundary algebra is the affine Kac--Moody vertex algebra
$A_\partial = V_k(\mathfrak{sl}_2)$ at level $k\neq -2$
(Theorem~\ref{thm:affine-half-space-bv}). Its generators are the
currents $J^a(z)$, $a\in\{e,f,h\}$, with OPE
\[
J^a(z)\,J^b(w)
\;\sim\;
\frac{k\,\kappa^{ab}}{(z-w)^2}
\;+\;
\frac{f^{ab}{}_c\,J^c(w)}{z-w}.
\]
The chiral Koszul dual (Vol~I, Theorem~A applied to the affine
lineage) is:
\begin{equation}\label{eq:su2-koszul-dual}
V_k(\mathfrak{sl}_2)^!
\;\simeq\;
V_{-k-4}(\mathfrak{sl}_2).
\end{equation}
This is the Feigin--Frenkel dual: the involution $k\mapsto -k-2h^\vee$
with $h^\vee = 2$ for $\mathfrak{sl}_2$. The Koszul duality is
\emph{same-family}: both sides are affine $\mathfrak{sl}_2$ algebras,
at dual levels.

The line-operator algebra, by the spectral braiding theorem
(Theorem~\ref{thm:spectral-braiding-core}) and the dg-shifted
factorization bridge
(Theorem~\ref{thm:complete-strictification}) in the proved
$\mathfrak{sl}_2$ path sector, is the strict Yangian
$Y_\hbar(\mathfrak{sl}_2)$; on the evaluation locus,
the line-operator category is
$\operatorname{Rep}_q(\mathfrak{sl}_2)$ with
$q = \exp(\pi i/(k+2))$ (from the log-HT monodromy identification
of Theorem~\ref{thm:affine-monodromy-identification}).

\medskip
\noindent\textbf{Step 2: the modular characteristic $\kappa$ and
the anomaly class.}
The modular characteristic of $V_k(\mathfrak{sl}_2)$ is computed from
the cyclic pairing on the bar complex (Vol~I, Theorem~D). For an
affine Kac--Moody algebra $V_k(\mathfrak{g})$, the modular
characteristic is
\begin{equation}\label{eq:km-kappa}
\kappaChHodge(V_k(\mathfrak{g}))
\;=\;
\frac{\dim\mathfrak{g}\cdot(k+h^\vee)}{2h^\vee}.
\end{equation}
For $\mathfrak{g}=\mathfrak{sl}_2$: $\dim\mathfrak{sl}_2=3$,
$h^\vee=2$, and
\begin{equation}\label{eq:su2-kappa}
\kappaChHodge(V_k(\mathfrak{sl}_2))
\;=\;
\frac{3(k+2)}{4}.
\end{equation}
This equals $\dim(\mathfrak{sl}_2)\cdot(k+h^\vee)/(2h^\vee)
= 3(k+2)/4$. Key properties:
\begin{enumerate}[label=\textup{(\roman*)}]
\item $\kappa = 0$ at $k=-2$ (critical level); $\kappa = 3/2$ at
 $k=0$, where the central double pole vanishes but the
 non-abelian simple-pole current algebra remains;
\item $\kappaChHodge(V_k) + \kappaChHodge(V_{-k-4}) = 3(k+2)/4 + 3(-k-2)/4 = 0$
 for all~$k$: complementarity is \emph{anti-symmetric} for
 Kac--Moody algebras (this anti-symmetry is special
 to KM, not universal);
\item $\kappa\neq 0$ for all $k\neq -2$.
\end{enumerate}

The anomaly class is
\begin{equation}\label{eq:su2-anomaly-class}
\Theta
\;:=\;
\kappaChHodge(V_k(\mathfrak{sl}_2))\cdot\omega_1
\;=\;
\frac{3(k+2)}{4}\cdot\omega_1
\;\in\;
Z^2(B)\cap Z(B),
\end{equation}
where $B = \bar B^{\mathrm{ch}}(V_k(\mathfrak{sl}_2))$ is the bar
dg~algebra and $\omega_1\in H^2(\overline{\mathcal{M}}_{1,1})$ is
the genus-$1$ tautological class ($\Theta$ is a \emph{class},
linear in~$\kappa$, not a scalar).

\medskip
\noindent\textbf{Step 3: the finite-dimensional shadow and the bar input.}
In the abelian case, the input algebra was $B=\Bbbk[x]$
(Chevalley--Eilenberg cochains of the abelian Lie algebra). For
\(\mathfrak{sl}_2\), the finite-dimensional shadow is the
Chevalley--Eilenberg cochain algebra
\begin{equation}\label{eq:su2-ce-input}
B
\;=\;
C^\bullet(\mathfrak{sl}_2)
\;=\;
\Lambda(\xi^e,\xi^f,\xi^h)
\end{equation}
with $|\xi^a|=1$ and CE differential
$d_{\mathrm{CE}}\xi^c = -\tfrac{1}{2}f^{ab}{}_c\,\xi^a\xi^b$.
Its cohomology is
\[
H^\bullet(\mathfrak{sl}_2;\Bbbk)
\cong \Bbbk\oplus \Bbbk\cdot\omega_3,
\]
matching the rational cohomology of \(SU(2)\simeq S^3\).

For the anomaly-completed theory, however, the relevant algebra is
not $C^\bullet(\mathfrak{sl}_2)$ itself but the bar dg~algebra
$B = \bar B^{\mathrm{ch}}(V_k(\mathfrak{sl}_2))$, which is an
infinite-dimensional dg~algebra controlling the full chiral OPE data.
The CE complex $C^\bullet(\mathfrak{sl}_2)$ appears as a
\emph{finite-dimensional shadow}: the genus-$0$ part of~$B$
receives a quasi-isomorphism from $C^\bullet(\mathfrak{sl}_2)$ on
the Koszul locus (one-loop exactness,
Theorem~\ref{thm:one-loop-koszul}).

\medskip
\noindent\textbf{Step 4: the transgression algebra.}
By Definition~\ref{def:tholog-transgression-algebra}, the
transgression algebra is
\begin{equation}\label{eq:su2-transgression}
B_\Theta
\;=\;
\frac{B*\Bbbk\langle\eta\rangle}
{\langle\eta b-(-1)^{|b|}b\eta\rangle},
\qquad
|\eta|=1,\qquad d\eta=\Theta,
\end{equation}
where the differential extends that of~$B$. The transgression
algebra \(B_\Theta\) is the universal neutralization of the
genus-\(1\) class on the affine bar model. Through the Koszul
comparison \(V_k(\mathfrak{sl}_2)^!\simeq
V_{-k-4}(\mathfrak{sl}_2)\), a \(B_\Theta\)-module \(M\) carries a
\(V_k(\mathfrak{sl}_2)\)-side bar action and a chain homotopy
\(h_M=\eta_M\) neutralizing the transported curvature action.

The secondary anomaly is
\begin{equation}\label{eq:su2-secondary}
u \;:=\; \eta^2 \;\in\; Z^2(B_\Theta)\cap Z(B_\Theta).
\end{equation}
By Proposition~\ref{prop:tholog-secondary-central-class}, $u$ is
closed and central. Unlike the abelian case where
$H^\bullet(B_\Theta)\cong \Bbbk[u]$, the cohomology of the
non-abelian transgression algebra is richer. By
Theorem~\ref{thm:tholog-hypersurface}, the long exact sequence
\[
\cdots\to
H^n(B)[u]\to H^n(B_\Theta)\to H^{n-1}(B)[u]
\xrightarrow{\cdot[\Theta]}
H^{n+1}(B)[u]\to\cdots
\]
computes $H^\bullet(B_\Theta)$ from \(H^\bullet(B)\) and the
\(\Theta\)-multiplication map. On the finite-dimensional CE shadow
\(C^\bullet(\mathfrak{sl}_2)\), \(H^2(\mathfrak{sl}_2;\Bbbk)=0\);
therefore the genus-\(1\) anomaly is invisible in that shadow. The
hypersurface quotient
\[
H^\bullet(B_\Theta)\cong H^\bullet(B)[u]/([\Theta])
\]
is asserted only for the full bar model, and only under the
injectivity hypothesis in Theorem~\ref{thm:tholog-hypersurface}.

\medskip
\noindent\textbf{Step 5: the anomaly $2$-cocycle as a Hochschild
class.}
The anomaly class $\Theta$ lives in $Z^2(B)\cap Z(B)$ and descends,
under the boundary reconstruction
(Proposition~\ref{prop:tholog-boundary-reconstruction}), to a
Hochschild class
\[
\Theta_\partial
\;=\;
\beta_{\mathrm{der}}(\Theta)
\;\in\;
HH^2(V_k(\mathfrak{sl}_2)).
\]
The Hochschild cohomology $HH^2(V_k(\mathfrak{sl}_2))$ controls
deformations of the VOA structure. The class $\Theta_\partial$ is
the genus-$1$ curvature obstruction, and its image under the
bulk--boundary quasi-isomorphism
$B\simeq HH^\bullet(V_k(\mathfrak{sl}_2))$ recovers~$\Theta$.

The key separation is that $\Theta_\partial$ is \emph{not}
the standard quantum group deformation cocycle (which lives in
$H^2(\mathfrak{g},\mathfrak{g})$ and controls deformations of
the \emph{Lie algebra structure}), but rather the \emph{genus-$1$
curvature class} (which lives in
$HH^2(V_k(\mathfrak{sl}_2))$ and controls deformations of the
\emph{vertex algebra structure} over the genus-$1$ moduli problem).
An evaluation functor from the Yangian side may compare the two
classes; the comparison obstruction is the difference between the
image of \(\Theta_\partial\) on evaluation modules and the standard
quantum-group deformation cocycle.

\medskip
\noindent\textbf{Step 6: the anomalous Steinberg for $SU(2)$ CS.}
Since \(\mathbb C\times\mathbb R\) is contractible, the derived
stack of flat \(SL_2\)-bundles is \(BSL_2\).  The non-abelian
field-theoretic local model used for the Steinberg comparison is the
shifted cotangent formal neighbourhood
\[
\Mvac^{\mathrm{for}}
\;=\;
T^*[-2]BSL_2
\;\simeq\;
[\mathfrak{sl}_2^*/SL_2],
\]
identified with \([\mathfrak{sl}_2/SL_2]\) after choosing the
Killing form.  The anomalous model carries a degree \(-2\)
pre-symplectic form \(\omega\) whose failure of closedness is the
't~Hooft anomaly:
\begin{equation}\label{eq:su2-anomaly-3form}
d\omega \;=\; H_3 \;=\; k\cdot c_3,
\end{equation}
where $c_3 = \tfrac{1}{6}\operatorname{Tr}(\theta\wedge[\theta,\theta])$
is the Cartan $3$-form on $SL_2$ (the generator of
$H^3(SU(2);\mathbb Z)\cong\mathbb Z$), pulled back to the
classifying stack, and $k$ is the Chern--Simons level.

The boundary Lagrangian is
\(\Lb = BSL_2 \hookrightarrow T^*[-2]BSL_2\)
(the Dirichlet/Nahm boundary condition, in the language of
Gaiotto--Witten), and the Steinberg object is
\[
\Sb
\;=\;
\Lb\times_{\Mvac^{\mathrm{for}}}\Lb,
\]
the derived self-intersection of the boundary Lagrangian in the
formal vacuum stack.

By Construction~\ref{constr:anomalous-steinberg}, the anomalous
Steinberg carries a \emph{curved} $(-1)$-shifted structure. The
curvature on the convolution algebra is:
\begin{equation}\label{eq:su2-curved-convolution}
d^2_{\mathrm{conv}}
\;=\;
[H_3,\;\cdot\;]
\;=\;
[k\cdot c_3,\;\cdot\;]
\;\neq\; 0.
\end{equation}
This is the non-abelian avatar of the curved bar differential
$d^2 = \kappa\cdot\omega_1$ from Vol~I: the Cartan $3$-form
$c_3$ plays the role of the genus-$1$ tautological class, and the
level $k$ plays the role of the modular characteristic~$\kappa$.

The convolution algebra on $\Sb$ is a curved dg~algebra: it is the
algebra of boundary operators with a non-vanishing curvature
proportional to $k\cdot c_3$. Modules over this curved algebra are
the \emph{anomaly-obstructed} boundary conditions of the $SU(2)$ CS
theory.

\medskip
\noindent\textbf{Step 7: anomaly completion, the extended stack
and the gerbe trivialization.}
Anomaly completion (Proposition~\ref{prop:anomaly-completion-trivialization})
uses an extended stack \(\widetilde\Mvac\to\Mvac^{\mathrm{for}}\)
on which the pullback of \(H_3\) is exact. For \(SU(2)\) CS this is
the trivialization problem for the Chern--Simons \(2\)-gerbe:
the \(3\)-form \(k\cdot c_3\) is the gerbe curvature, and the
extended stack carries a \(2\)-form potential \(\beta\) with
\(d\beta=k\cdot c_3\).

Concretely, for $SU(2)$:
\begin{enumerate}[label=\textup{(\roman*)}]
\item The Cartan $3$-form $c_3$ generates
 $H^3(SU(2);\mathbb Z)\cong\mathbb Z$. At level~$k$,
 $k\cdot c_3$ is $k$ times the generator.
\item One model of the trivialization is a lift through the
 string \(2\)-group, equivalently a trivialization of the
 Chern--Simons gerbe class. In the algebraic model this is encoded by
 a central gerbe extension
 \[
 1 \;\to\; B\mathbb G_m \;\to\;
 \widetilde\Mvac \;\to\; \Mvac^{\mathrm{for}} \;\to\; 1,
 \]
 where the $B\mathbb G_m$-fiber is the gerbe data trivializing
 $k\cdot c_3$.
\item At the level of cochain algebras, the transgression comparison is
 \begin{equation}\label{eq:su2-transgression-match}
 C^\bullet(\widetilde\Mvac)
 \;\simeq\;
 C^\bullet(\Mvac^{\mathrm{for}})_\Theta
 \;=\;
 B_\Theta.
 \end{equation}
 This isomorphism is an input comparison: the transgression generator
 \(\eta\) maps to the cochain-level \(2\)-form potential, and
 \(d\eta=\Theta\) maps to \(d\beta=k\cdot c_3\) under the chosen
 normalization.
\end{enumerate}

The extended Steinberg is
\begin{equation}\label{eq:su2-extended-steinberg}
\widetilde\Sb
\;=\;
\widetilde\Lb
\;\times_{\widetilde\Mvac}\;
\widetilde\Lb,
\end{equation}
where $\widetilde\Lb$ is the lift of the boundary Lagrangian to the
extended stack. On $\widetilde\Sb$, the curvature vanishes:
\[
d^2_{\mathrm{conv}}\big|_{\widetilde\Sb} \;=\; 0.
\]
The convolution algebra on $\widetilde\Sb$ is an \emph{honest}
(uncurved) dg~algebra, and its modules are the
\emph{anomaly-completed} boundary conditions.

\medskip
\noindent\textbf{Step 8: the comparison, transgression algebra
vs.\ uncurved convolution.}
If the genus-$0$ Steinberg convolution identifies with the affine
bar dg~algebra and the gerbe trivialization identifies with the
cochain transgression of Step~7, the comparison is
\begin{equation}\label{eq:su2-comparison}
\text{convolution algebra on }\widetilde\Sb
\;\simeq\;
B_\Theta
\;=\;
(V_k(\mathfrak{sl}_2)^!)_\Theta,
\end{equation}
where the left side is the anomaly-completed convolution algebra
(geometric) and the right side is the transgression algebra
(algebraic). The comparison has three components:
\begin{enumerate}[label=\textup{(\alph*)}]
\item \emph{Genus-$0$:} the uncompleted convolution algebra on $\Sb$
 recovers the bar dg~algebra $B = \bar B^{\mathrm{ch}}(V_k(\mathfrak{sl}_2))$,
 which controls twisting morphisms between $V_k(\mathfrak{sl}_2)$ and
 $V_{-k-4}(\mathfrak{sl}_2)$ (the bar complex classifies
 twisting morphisms, not bulk observables);
\item \emph{Anomaly completion:} the passage from $\Sb$ to
 $\widetilde\Sb$ is, at the cochain level, precisely the passage
 from $B$ to $B_\Theta$ (the transgression, Step~7(iii));
\item \emph{$(-1)$-shifted symplectic structure:} the honest
 $(-1)$-shifted symplectic form on $\widetilde\Sb$ descends from
 Vol~I's complementarity theorem (Theorem~C): the
 decomposition of genus-$g$ obstructions into complementary
 Lagrangians lifts to the extended locus, where the shifted
 symplectic form is uncurved.
\end{enumerate}

\medskip
\noindent\textbf{Step 9: the secondary anomaly and the quantum
Casimir.}
The secondary anomaly $u = \eta^2\in Z^2(B_\Theta)\cap Z(B_\Theta)$
acquires a concrete identification in the $SU(2)$ theory.

On the Yangian side ($Y_\hbar(\mathfrak{sl}_2)$ as line-operator
algebra), the anomaly-completed Yangian
$Y_{\hbar,\Theta} := (Y_\hbar(\mathfrak{sl}_2))_\Theta$
(Theorem~\ref{thm:tholog-dg-shifted-yangian-completion})
carries $u = \eta^2$ as a central element. On the evaluation
locus $\operatorname{Rep}_q(\mathfrak{sl}_2)$, fix the quantum
Casimir eigenvalue
\[
C_q(\lambda)=q^{\lambda+1}+q^{-\lambda-1}-(q+q^{-1})
\]
on the highest-weight module \(V_\lambda\).  The normalization
defect of the Casimir comparison is
\begin{equation}\label{eq:su2-secondary-casimir}
\mathcal E_\lambda
\;:=\;
u_{V_\lambda}
-
\frac{\kappa}{k+2}
\cdot
C_q(\lambda),
\end{equation}
where the prefactor \(\kappa/(k+2)\) is the chosen normalization
between the chiral modular characteristic and the quantum group
parameter.  The equality with the quantum Casimir is the exact
condition \(\mathcal E_\lambda=0\) for all evaluation modules.  This
is the non-abelian analogue of the abelian relation \(d\eta=kx\):
the scalar curvature \(kx\) is replaced by the representation-dependent
Casimir eigenvalue.

\medskip
\noindent\textbf{Step 10: genus-Clifford completion.}
The genus-$g$ Clifford completion
(Definition~\ref{def:tholog-genus-clifford})
\[
\mathfrak G_g(B_\Theta)
\;=\;
\frac{B_\Theta\langle\alpha_1,\beta_1,\dots,\alpha_g,\beta_g\rangle}
{\langle\alpha_i^2,\;\beta_j^2,\;
\alpha_i\beta_j+\beta_j\alpha_i-\delta_{ij}\eta^2\rangle}
\]
is the genus-\(g\) handle extension of the transgression algebra.

\emph{Away from $u=0$.}
When $\eta^2$ acts invertibly, Morita triviality
(Theorem~\ref{thm:tholog-morita-triviality}) gives
\[
\mathfrak G_g(B_\Theta)[\eta^{-2}]
\;\cong\;
\operatorname{Mat}_{2^g}\bigl(B_\Theta[\eta^{-2}]\bigr).
\]
Higher genus is Morita-invisible after this localization. This is a
statement about the localized genus-Clifford completion, not about the
entire \(V_k(\mathfrak{sl}_2)\)-module category.

\emph{The strict locus $u=0$.}
Setting $\eta^2=0$ (Theorem~\ref{thm:tholog-exterior-degeneration}),
\[
\mathfrak G_g(B_\Theta)/(\eta^2)
\;\cong\;
(B_\Theta/(\eta^2))
\otimes
\Lambda(\alpha_1,\beta_1,\dots,\alpha_g,\beta_g),
\]
so the handle algebra is exterior.  A further replacement of
\(B_\Theta/(\eta^2)\) by a quotient of \(B\) requires the
hypersurface injectivity hypothesis of
Theorem~\ref{thm:tholog-hypersurface}.

\emph{The dichotomy at critical level.}
At $k=-2$ (critical level): $\kappa = 3(-2+2)/4 = 0$, the
Sugawara construction is undefined, so the affine topologisation
arguments using the Sugawara stress tensor do not apply.  The scalar
formula gives zero curvature, but the critical-level center requires a
separate analysis. At \(k=0\), the central double pole vanishes while
the non-abelian simple-pole current algebra remains; the modular
characteristic is \(\kappa=3/2\).

\medskip
\noindent\textbf{Summary.}
For the \(SU(2)\) comparison at non-critical level \(k\neq -2\):
\begin{center}
\renewcommand{\arraystretch}{1.2}
\begin{tabular}{ll}
Boundary algebra & $V_k(\mathfrak{sl}_2)$ \\
Koszul dual & $V_{-k-4}(\mathfrak{sl}_2)$ (Feigin--Frenkel) \\
Modular characteristic &
 $\kappa = 3(k+2)/4$ \\
Complementarity & $\kappa + \kappa' = 0$ (KM anti-symmetry) \\
Anomaly class &
 $\Theta = \kappa\cdot\omega_1 \in Z^2(B)\cap Z(B)$ \\
Formal vacuum model & $\Mvac^{\mathrm{for}}=T^*[-2]BSL_2$ \\
Anomaly $3$-form & $H_3 = k\cdot c_3$ (Cartan $3$-form) \\
Anomalous Steinberg &
 $\Sb = \Lb\times_{\Mvac^{\mathrm{for}}}\Lb$, curved $(-1)$-shifted \\
Extended stack &
 $\widetilde\Mvac$: gerbe trivialization of $k\cdot c_3$ \\
Transgression algebra &
 $B_\Theta = B\langle\eta\mid d\eta=\Theta\rangle$ \\
 & compared with convolution on $\widetilde\Sb$ \\
Secondary anomaly & $u=\eta^2$; on eval.\ locus
 defect $\mathcal E_\lambda$ measures Casimir normalization \\
Genus-$g$ completion &
 $\mathfrak G_g(B_\Theta)$: Clifford on $2g$ generators \\
Morita triviality &
 $\mathfrak G_g[\eta^{-2}]
 \cong\operatorname{Mat}_{2^g}(B_\Theta[\eta^{-2}])$ \\
Strict locus &
 $\mathfrak G_g/(\eta^2)
 \cong (B_\Theta/(\eta^2))\otimes\Lambda^{2g}$
\end{tabular}
\end{center}

\emph{What is new relative to the abelian case.}
Three features are genuinely new:
\begin{enumerate}[label=\textup{(\arabic*)}]
\item The input algebra $B$ is non-commutative (it is the bar
 dg~algebra of a non-abelian VOA), so the transgression algebra
 $B_\Theta$ is a non-commutative Ore extension, not a polynomial
 ring.
\item The anomalous Steinberg $\Sb$ is a nontrivial derived
 intersection (the derived self-intersection of
 \(BSL_2\) in \(T^*[-2]BSL_2\)), not a
 product of points. The curved $(-1)$-shifted structure encodes
 the full non-abelian Chern--Simons $3$-form, not just a level
 times a curvature.
\item The secondary anomaly $u=\eta^2$ has a quantum-Casimir
 comparison on evaluation modules; the defect
 \(\mathcal E_\lambda\) records the normalization error. In the
 abelian case the curvature is the scalar \(kx\); in the
 non-abelian case the comparison varies across the weight lattice.
\end{enumerate}
\end{construction}

\begin{remark}[The non-abelian gerbe as string structure]
\label{rem:su2-string-structure}
\index{string structure!anomaly completion}
The passage from $\Mvac^{\mathrm{for}}$ to $\widetilde\Mvac$ in Step~7 is
an instance of the general principle that anomaly completion
trivializes a \emph{higher gerbe}. For $SU(2)$:
\begin{itemize}
\item $H^3(SU(2);\mathbb Z)\cong\mathbb Z$ is generated by $c_3$;
\item the level-$k$ Chern--Simons theory picks out $k\cdot c_3$;
\item trivialization of this class requires a \emph{string structure}
 (a trivialization of the first fractional Pontryagin class, or
 equivalently a lift through the $3$-connected cover of $SU(2)$).
\end{itemize}
The transgression generator $\eta$ is the cochain-level witness of
this string trivialization: $d\eta = \Theta$ states that the anomaly
class becomes exact on the string-extended stack. The secondary
anomaly \(u=\eta^2\) is the quadratic class of the chosen
trivialization.  Its vanishing is the strict-locus condition for the
genus-Clifford algebra.

For higher-rank simple simply connected compact groups \(G\), the
same construction applies after choosing the corresponding gerbe
trivialization:
$c_3$ generates $H^3(G)$, the CS level selects $k\cdot c_3$, and
the anomaly-completed transgression algebra trivializes it at the
cochain level. The rank of the anomalous Steinberg and the
complexity of the quantum Casimir grow with
$\dim\mathfrak{g}$, but the structural framework
(transgression, secondary anomaly, genus-Clifford, Morita
dichotomy) is the same.
\end{remark}

\begin{remark}[Comparison with the $\mathfrak{sl}_3$ DS chain]
\label{rem:su2-vs-sl3-chain}
The $SU(2)$ computation above sits at level~$0$ of the
$\mathfrak{sl}_2$ DS chain: since $\mathfrak{sl}_2$ has only two
nilpotent orbits ($\{0\}$ and the regular orbit), the full chain is
\[
\widehat{\mathfrak{sl}}_{2,k}
\;\xrightarrow{\;\mathrm{DS}_{\mathrm{prin}}\;}
\mathrm{Vir}_{c(k)},
\qquad
c(k) = 1-\frac{6(k+1)^2}{k+2}.
\]
Under the hypotheses of the DS comparison criterion
(Theorem~\ref{thm:ds-anomaly-functoriality}), the anomaly-completed
datum has the following two-vertex comparison:
\begin{center}
\renewcommand{\arraystretch}{1.2}
\begin{tabular}{lcc}
& \textbf{Level~$0$} & \textbf{Level~$1$} \\
\hline
Algebra & $V_k(\mathfrak{sl}_2)$ & $\mathrm{Vir}_{c(k)}$ \\[2pt]
Koszul dual & $V_{-k-4}(\mathfrak{sl}_2)$ & $\mathrm{Vir}_{26-c(k)}$ \\[2pt]
$\kappa$ & $3(k+2)/4$ & $c/2$ \\[2pt]
$\kappa+\kappa'$ & $0$ & $13$ \\[2pt]
Self-dual point & (none; anti-symmetric) & $c=13$ \\[2pt]
Anomalous Steinberg & $\Sb$ with $H_3=k\cdot c_3$ & $\Sb'$ with
 $H_3'=\kappa_{\mathrm{Vir}}\cdot\omega_1$
\end{tabular}
\end{center}
The DS reduction compares the non-abelian Cartan $3$-form
$k\cdot c_3$ into the tautological genus-$1$ class
$\kappa_{\mathrm{Vir}}\cdot\omega_1$ for the Virasoro algebra:
the ghost system absorbs the non-abelian structure of
$\mathfrak{sl}_2$, leaving only the genus-$1$ modular
characteristic as the residual anomaly when the curvature transform
\eqref{eq:curvature-ds-transform} holds.
\end{remark}

\section{Mukai-doubling conductor $K^{\kappaChHodge} = 8$
on the Lorentzian-lattice $\mathcal{B}$-family
and the Theorem~C enlargement}
\label{sec:ac-comparison-supplement}

\begin{remark}[Mukai doubling and Bruinier reciprocity]
\label{rem:ac-comparison-mukai-bruinier}
For anomaly completion, the Mukai-doubling Koszul conductor $K^{\kappaChHodge}(\mathcal{H}_{\mathrm{Muk}}(K3)) = 2\,c_+(\mathrm{Mukai}(K3)) = 8$ identifies (via Bruinier $2002$ Proposition~$5.1$ Heegner Chern-class reciprocity) with: (a) order of monodromy of $\mathcal{L}^{\Delta_5}$ around Humbert divisor $H_1 \subset \overline{\mathcal{A}_2}$; (b) Lusztig specialisation $\ell = 8$ at $\zeta^8 = 1$, $\hbar^2 = -1/8$. Universal identity $\hbar^2 \cdot K^{\kappaChHodge} = -1$ on the Lorentzian-lattice-parametric $\mathcal{B}$-family. Theorem~C enlarges to $\{0, 8, 13, 250/3, 98/3\}$.
\end{remark}
